\newpage
\ssscp{\tsc{South Babar}}{04-06-05-02}
A \tsc{South Babar} node containing all remaining \tsc{Babar}
languages is well justified on the basis of five changes,
given in \qf{04-28} below.
Of these changes, *s > **ʈ > \it{t} is the only exclusively shared
innovation.

{\wg\st{6pt}\tblr{>{\hspace{-\tabcolsep}}l<{\hspace{-\tabcolsep}}>{\hspace{-\tabcolsep}\enspace}lllll}
\tbx{04-28}	&\mc{5}{l}{\tsc{South Babar} innovations}\\
						&a.	&*k				&>&{\hr\hr}{\0}	&\\
						&b.	&*t				&>&**k	&\\
						&c.	&*s > **ʈ	&>&**t	&\\
						&d.	&*a/*ə		&>&{\hr\hr}{\0}	&/σ\und{σ}{\#}\\
						&e.	&*ŋ				&>&**n	&\\
\end{tabular}\mbox{}\\}

The change *k > {\0} began in Proto-Babar,
but there are some indications that it was retained as **ʔ in
some Proto-South Babar words (see \trf{04-57}).
Loss of non-high vowels also occurs in Dawera
and Dai (see \trf{04-51}),
and *ŋ > \it{n} is common in all \tsc{Timor-Babar} languages.
Similarly, while North Babar has some cases of *t > \it{k},
it is a conditioned change in that language, while it is unconditioned in \tsc{South Babar}.

Examples of *s > **ʈ > \it{t} in \tsc{South Babar} have already
been given in \trf{04-45} (\prf{tab:04-45}).
Examples of *t > \it{k} are given in \trf{04-54}.
Serili then has intervocalic **k > \it{k,
ʔ,} {\0}, Southeast Babar has non-initial **k > \it{x},
and Central Marsela has non-initial **k > \it{k}, {\0}.\footnote{
		*t was retained as \it{t} after a nasal,
		probably through *nt > **nd > **ʈ with subsequent regular **ʈ > \it{t}.
		This is seen in reflexes of *punti ‘banana’ in \trf{04-01}.}

\begin{table}\tcp{\tsc{South Babar} *t > \it{k}}{04-54}
\begin{adjustbox}{width=\textwidth}
	\begin{threeparttable}\st{2pt}
	\begin{tabular}{lll@{}l@{}l@{}l@{}l}\lsptoprule
local gl.	&	{\hr}star	&	{\hr}knee	&	{\hr}faeces	&	{\hr}vomit	&	{\hr}stone	&	{\hr}eye\\
PMP	&	\hr*bi\tbr{t}uqən	&	\hr*\tbr{t}uhud	&	\hr*\tbr{t}aqi	&	\hr*um-u\tbr{t}aq	&	\hr*ba\tbr{t}u	&	\hr*ma\tbr{t}a\su{Δ}\\
P.-Babar	&	**\tbr{t}un	&	**\tbr{t}ur	&	**\tbr{t}ei	&	**mu\tbr{t}a	&	**wa\tbr{t}u	&	**ma\tbr{t}a\\
P.-S. Babar	&	**\tbr{k}un	&	**\tbr{k}ur	&	**\tbr{k}ei	&	**mu\tbr{k}u	&	**wa\tbr{k}i	&	**mu\tbr{k}a\\ \midrule
C. Marsela	&	{\hr\hr}\tbr{k}un	&	{\hr\hr}\tbr{k}or-no	&	{\hr\hr}\tbr{k}e-ni-ei\su{†}	&	{\hr\hr}nɛ-mu\tbr{k}	&	{\hr\hr}way	&	{\hr\hr}mo-nei,\\
	&	&		&		&		&		&	{\hr\hr}mo\tbr{kk}o\\
W. Marsela	&	{\hr\hr}\tbr{k}un	&	{\hr\hr}\tbr{k}on-nɛ	&	{\hr\hr}\tbr{k}e	&	{\hr\hr}nɛ-mo\tbr{k}	&	{\hr\hr}wa\tbr{k}y	&	{\hr\hr}mu\tbr{k}\\
E. Marsela	&	{\hr\hr}\tbr{k}un	&	{\hr\hr}\tbr{k}or-l	&	{\hr\hr}\tbr{k}ei	&	{\hr\hr}mu\tbr{k}	&	{\hr\hr}wa\tbr{k}	&	{\hr\hr}mo\tbr{k}\\
Serili	&	{\hr\hr}\tbr{k}ul, ul-ɛ	&	{\hr\hr}\tbr{k}or	&	{\hr\hr}\tbr{k}e	&	\hr\cg{(l\<y\>enoo)}	&	{\hr\hr}ba\tbr{ʔ}ay-e\su{‡}	&	{\hr\hr}mɔɔ, mo\tbr{ʔ}o\\
SE. Babar	&	{\hr\hr}\tbr{k}ulamol	&	{\hr\hr}\tbr{k}or	&	{\hr\hr}\tbr{k}e	&	{\hr\hr}mo\tbr{x}	&	{\hr\hr}wa\tbr{x}ay\su{‡}	&	{\hr\hr}mo\tbr{x}\\
Emplawas	&	{\hr\hr}\tbr{k}ul-o	&	{\hr\hr}\tbr{k}or-wo	&	{\hr\hr}\tbr{k}ɛ-o	&	{\hr\hr}lɛ-mo\tbr{k}	&	{\hr\hr}wa\tbr{k}y-o	&	{\hr\hr}mu\tbr{k}-lɛ\\
Imroing	&	\hr\cg{(kowa{\tl}kowa)}	&	\hr\cg{(weker-mu)}	&	{\hr\hr}\tbr{k}e-mu-te	&	{\hr\hr}mo-mo\tbr{k}u	&	{\hr\hr}wa\tbr{k}y-e	&	{\hr\hr}mu\tbr{k}u-mu-te\\
Tela	&	{\hr\hr}\tbr{k}un-e	&	{\hr\hr}\tbr{k}or-ne	&	{\hr\hr}\tbr{k}e-nie\su{†}	&	{\hr\hr}na-mo\tbr{k}y	&	{\hr\hr}wa\tbr{k}-e	&	{\hr\hr}mu\tbr{k}-ni\\
		\lspbottomrule
	\end{tabular}
		\begin{tablenotes}\tnsize
			\item[†]	{The stems of Central Marsela \it{ke-ni-ei}
								and Tela \it{ke-nie} ‘faeces’ are \it{kei-} with metathesis of
								the final high vowel across a morpheme boundary.}
			\item[‡]	{Serili \it{baʔay-e}
								and Southeast Babar \it{waxay} ‘stone’ have diphthongisation
								of final \it{i}.
								%as described by \citet[96]{Steinhauer2009}.
								This is also seen in reflexes of *kutu > **oti ‘louse’ in \trf{04-56}.
								\citet[158]{vanDijk2021} gives Serili ‹baa› ‘stone’,
								which is probably phonetically [baʔa].}
			\item[Δ]	{C. Marsela \it{mo-nei} `eye' is from \citet[414]{Taber1993}
								and \it{mokko} from \citet[163]{vanDijk2021}.
								Serili \it{mɔɔ} [mɔˈɔ] is from \citet[414]{Taber1993}
								and \it{moʔo} ‹mo{\Q}o› is from \citet[163]{vanDijk2021}.}
		\end{tablenotes}
	\end{threeparttable}
\end{adjustbox}
\end{table}

Within \tsc{South Babar} we can further identify a \tsc{Nuclear
South Babar} node which includes all languages except Central Marsela.
This node is defined by word-initial *a > \it{u}.
This vowel also shifts after word-initial nasals,
with *a > \it{u} /{\#}N{\gap} in Emplawas,
Imroing, and Tela,\footnote{
Tela has sporadic *a > \it{o} /{\#}(N){\gap}
instead of expected *a > \it{u}.} 
and *a > \it{o} /{\#}N{\gap} in East Marsela,
Serili, and Southeast Babar.
Examples are given in \trf{04-55}.
Central Marsela has *a > \it{e} in the same environment.\footnote{
		There are two words in which Central Marsela has
		*a > \it{o} /{\#}m{\gap}.
		PMP *mata > **muka > \it{mo-nei},
		\it{mokko} ‘eye’ (\trf{04-54}),
		and *matay > **muki > \it{n-mokei} ‘die’.}
\citet{Zobel2018} proposes that these reflexes are all from intermediate **ə.
Thus, he proposes PMP *a > Proto-South Babar **ə /{\#}(N){\gap},
with **ə > \it{e} in Central Marsela,
and subsequent **ə > \it{u} and/or *ə > \it{o} /{\#}(N){\gap} in other languages.

\begin{table}\tcp{\tsc{Nuclear \tsc{South Babar}} *a > \it{u} /{\#}(N){\gap}}{04-55}
\begin{adjustbox}{width=\textwidth}
	\begin{threeparttable}\st{2pt}
	\begin{tabular}{lllllllll}\lsptoprule
local gl.	&	{\hr}\tsc{1sg}		&	{\hr}child					&	{\hr}ash						&	{\hr}fire						&	{\hr}bird							&	{\hr}name						&	{\hr}swim								&	{\hr}tongue\\
PMP				&	\hr*\tbr{a}ku			&	\hr*\tbr{a}nak			&	\hr*q\tbr{a}bu			&	\hr*h\tbr{a}puy			&	\hr*m\tbr{a}nuk				&	\hr*ŋ\tbr{a}jan			&	\hr*n\tbr{a}ŋuy					&	\\
P.Babar		&	**\tbr{a}u				&	**\tbr{a}na					&	**(li)-\tbr{a}wi		&	**\tbr{a}i					&	**m\tbr{a}nu					&	**ŋ\tbr{a}n					&	**wen\tbr{a}ŋui					&	**n\tbr{a}ma\su{†}\\
P.S. Babar&	**\tbr{a}i				&	**\tbr{a}n					&	**\tbr{a}wi					&	**\tbr{a}i					&	**m\tbr{a}ni(en)			&	**n\tbr{a}n					&	**-wen\tbr{a}ni					&	**n\tbr{a}m\\	\midrule
C. Marsela&	{\hr\hr}\tbb{ɛ}		&	{\hr\hr}\tbb{ɛ}n-ei	&	{\hr\hr}\tbb{ɛ}wy		&	{\hr\hr}\tbb{ɛ}y-ei	&	{\hr\hr}m\tbb{ɛ}nyɛn	&	{\hr\hr}n\tbb{ɛ}n		&	{\hr\hr}ne-wy\tbb{ɛ}n		&	{\hr\hr}n\tbb{ɛ}m-nei\\
W. Marsela&	{\hr\hr}\tbr{u}(i)&	{\hr\hr}\tbr{u}n		&	{\hr\hr}\tbr{u}wy		&	{\hr\hr}\tbr{u}y		&	{\hr\hr}m\tbr{u}ny		&	{\hr\hr}n\tbr{u}n		&	{\hr\hr}mo-wny\tbb{ɛ}n	&	{\hr\hr}n\tbr{u}m\\
E. Marsela&	{\hr\hr}\tbr{u}i	&	{\hr\hr}\tbr{u}n		&	{\hr\hr}\tbr{u}w		&	{\hr\hr}\tbr{u}i		&	{\hr\hr}m\tbr{o}nyen	&	{\hr\hr}n\tbr{o}n		&	{\hr\hr}mo-wn\tbr{o}n		&	{\hr\hr}n\tbr{o}m\\
Serili		&	{\hr\hr}\tbr{u}i	&	{\hr\hr}\tbr{u}n-e	&	{\hr\hr}\tbr{u}bi-e	&	{\hr\hr}\tbr{u}y-e	&	{\hr\hr}m\tbr{o}n-e		&	{\hr\hr}n\tbr{o}n		&	{\hr\hr}lyɛ-mo-n\tbr{o}n&	{\hr\hr}n\tbr{o}m\\
SE. Babar	&	{\hr\hr}\tbr{u}m	&	{\hr\hr}\tbr{u}n		&	{\hr\hr}\tbr{u}bi		&	{\hr\hr}\tbr{u}y		&	{\hr\hr}m\tbr{o}nyan	&	{\hr\hr}n\tbr{o}n		&	{\hr\hr}ke-ben\tbr{o}n	&	{\hr\hr}n\tbr{o}m\\
Emplawas	&	{\hr\hr}n\tbr{u}	&	{\hr\hr}\tbr{u}n-ɛ	&	{\hr\hr}\tbr{u}wy-o	&	{\hr\hr}\tbr{u}y-o	&	{\hr\hr}m\tbr{u}nyɛl-o&	{\hr\hr}n\tbr{u}n-o	&	{\hr\hr}lɛ-wn\tbr{u}n		&	{\hr\hr}n\tbr{u}p-ɛ\\
Imroing		&	{\hr\hr}\tbr{u}i	&	{\hr\hr}\tbr{u}n-emu&	{\hr\hr}\tbr{u}wi-e	&	{\hr\hr}\tbr{u}y-e	&	{\hr\hr}m\tbr{u}n-e		&	{\hr\hr}n\tbr{u}n-mu&	{\hr\hr}mo-l\tbr{u}li		&	{\hr\hr}n\tbr{u}m-mu\\
Tela			&	{\hr\hr}\tbr{u}i	&	{\hr\hr}\tbr{u}n-e	&	{\hr\hr}\tbr{o}bi-e	&	{\hr\hr}\tbr{u}y-e	&	{\hr\hr}m\tbr{u}n-e		&	{\hr\hr}n\tbr{u}n-mu&	{\hr\hr}ne-wn\tbr{u}ni	&	{\hr\hr}n\tbr{o}m-ne\\
		\lspbottomrule
	\end{tabular}
		\begin{tablenotes}\tnsize
			\item[†] {Proto-Babar **nama ‘tongue’ is cognate with Kisar \it{naman}, Luang \it{namə} ‘tongue’.}
		\end{tablenotes}
	\end{threeparttable}
\end{adjustbox}
\end{table}

Whatever the exact sequence of changes,
*a > \it{o} /{\#}N{\gap} provides phonological evidence for placing
East Marsela, Serili, and Southeast Babar together.
These languages also group together (with Central Marsela) on
the basis of lexical similarity in \citet[407f]{Taber1993}.

Given cross-linguistic patterns of sound change,
we might expect \it{o} as intermediate between *a and \it{u}.
However, in certain cases PMP *u > Proto-South Babar **o without
merging with PMP *a.\footnote{
		The usual environment for *u >
		Proto-South Babar **o is before open syllables with a non-high
		vowel (/{\gap}(C)V\tsc{[-high]}{\#}),
		as well as before **k < *t \citep[103]{Steinhauer2009},
		but exceptions occur.}
This means that PMP *a cannot have
first become Proto-South Babar **o before shifting to \it{u}.\footnote{
		The discussion and analysis in this section may need slight
		revision in light of Edwards' recent fieldwork, as both Southeast Babar
		and Serili have distinct mid-high and mid-low vowels.
		They have *a > \it{ɔ} /N{\gap}, but *u > \it{o} in the examples
		in \trf{04-56}.}
Examples are given in \trf{04-56} to illustrate.
Thus, *a > \it{u} /{\#}(N){\gap} cannot have gone through intermediate **o.
This does not preclude the possibility that *a > \it{u} passed
through an alternate intermediate stage,
such as a centralised vowel **ə,
as proposed by \citet{Zobel2018}.

\begin{table}\tcp{\tsc{South Babar} **o}{04-56}
\begin{adjustbox}{width=\textwidth}
	\begin{threeparttable}\st{2pt}
	\begin{tabular}{llll@{}l@{}ll}\lsptoprule
local gl.	&	bamboo	&	\tsc{2sg}	&	head	&	louse	&	wound	&	fingernail\\
PMP	&	*q\tbr{au}R	&	*k\tbr{ahu}	&	*h\tbr{u}tək\su{†}	&	*k\tbr{u}tu	&	*(ma)-n\tbr{u}ka	&	\\
P.-Babar	&	**\tbr{o}r	&	**\tbr{o}u	&	**k\tbr{u}tu	&	**k\tbr{u}tu	&	**(m)n\tbr{u}	&	\\
P.-S. Babar	&	**\tbr{o}r	&	**\tbr{o}wu	&	**\tbr{o}ka	&	**\tbr{o}ki	&	**mn\tbr{o}	&	**\tbr{o}m\\	\midrule
C. Marsela	&	{\hr}\tbr{o}r	&	{\hr}\tbr{o}u	&	{\hr}\tbr{o}k-e	&	{\hr}\tbr{o}i-ei	&	{\hr}n\tbr{o}-ei	&	{\hr}\tbr{o}m-ei\\
W. Marsela	&	{\hr}\tbr{o}r	&	{\hr}\tbr{o}w	&	{\hr}\tbr{o}k	&	{\hr}\tbr{o}ky	&	{\hr}n\tbr{o}{\tl}no	&	{\hr}\tbr{o}m\\
E. Marsela	&	{\hr}\tbr{o}r	&	{\hr}\tbr{o}wi	&	{\hr}\tbr{o}k-e	&	{\hr}\tbr{o}k	&	{\hr}l\tbr{o}	&	{\hr}\tbr{o}m\\
Serili	&	{\hr}\tbr{ɔ}r	&	{\hr}\tbr{o}bi	&	{\hr}\tbr{o}o, h\tbr{o}ʔo\su{‡}	&	{\hr}h\tbr{o}ʔoy-e\su{‡}	&	{\hr}l\tbr{o}	&	{\hr}\tbr{o}m\\
SE. Babar	&	\hr\cg{(bex-e)}	&	{\hr}\tbr{o}b	&	{\hr}\tbr{o}x	&	{\hr}\tbr{o}xoy	&	{\hr}ml\tbr{o}	&	{\hr}\tbr{o}um\\
Emplawas	&	{\hr}\tbr{o}r-o	&	{\hr}\tbr{o}	&	{\hr}\tbr{o}k-ɛ	&	{\hr}\tbr{o}kɛ	&	{\hr}lɛ-l\tbr{o}	&	{\hr}\tbr{o}m-o\\
Imroing	&	{\hr}l\tbr{o}r-le	&	{\hr}\tbr{o}wu	&	{\hr}\tbr{o}ka-mu	&	{\hr}\tbr{o}ky-e	&	{\hr}le-l\tbr{o}	&	\hr\cg{(ri-ɛ muk-mu-te)}\\
Tela	&	\hr\cg{(pɛkk-e)}	&	{\hr}\tbr{o}wi	&	{\hr}\tbr{o}k-e	&	{\hr}\tbr{o}k-e	&	{\hr}ne-mn\tbr{o}	&	{\hr}\tbr{o}m-e\\
		\lspbottomrule
	\end{tabular}
		\begin{tablenotes}\tnsize
			\item[†]	{PMP *hutək is reconstructed with the meaning ‘brain’.}
			\item[‡]	{Serili \it{hoʔo} ‘head’
								and \it{hoʔoy-e} ‘louse’ have insertion of initial
								\it{h}, before a medial glottal stop.
								A similar process occurs in some Rote languages
								and East Tetun (see *hapuy ‘fire’ in \trf{04-01} on \prf{tab:04-01}).
								Serili \it{oo} [oˈo] ‘head’ is from \citet[423]{Taber1993}
								and \it{hoʔo} from \citet[162]{vanDijk2021}.}
		\end{tablenotes}
	\end{threeparttable}
\end{adjustbox}
\end{table}

There are a several words which have initial \it{a} in nearly
all \tsc{South Babar} languages,
as shown in \trf{04-57}.
In two cases (**halan ‘road’
and **har ‘betel pepper’),
Serili has initial \it{h}
and in one case (**yai `boat') external cognates have initial \it{y}.
For these words we can thus posit that these consonants were
still present when *a > \it{u} /{\#}(N){\gap} was taking place.
Thus, the vowel was “protected” by the preceding consonant.
The remaining examples with initial *a = \it{a} are reflexes
of words with initial PMP *k.
For these words we reconstruct an initial glottal stop at the
level of Proto-South Babar which was lost after initial *a > \it{u}.
A similar approach is taken by \citet{Zobel2018}.

\begin{table}\tcp{\tsc{South Babar} *a = \it{a}}{04-57}
%\begin{adjustbox}{width=\textwidth}
	\begin{threeparttable}\st{2pt}
	\begin{tabular}{lllllll}\lsptoprule
local gl.	&	{\hr}road	&	betel pepper	&	{\hr}elder sibling	&	{\hr}\tsc{1pl.excl}	&	{\hr}stomach	&	boat\\
PMP	&	\hr*z\tbr{a}lan	&		&	\hr*k\tbr{a}ka	&	\hr*k\tbr{a}mi	&	\hr*k\tbr{a}mbu	&	\\
P.-Babar	&	**h\tbr{a}lan	&	**h\tbr{a}r\su{†}	&	**k\tbr{a}ka	&	**k\tbr{a}mi	&	**ʔ\tbr{a}pun	&	**y\tbr{a}i\\
P.-S. Babar	&	**h\tbr{a}ln	&	**h\tbr{a}r	&	**ʔ\tbr{a}ka	&	**ʔ\tbr{a}mi	&	**ʔ\tbr{a}p(e)n	&	**y\tbr{a}i\\ \midrule
C. Marsela	&	{\hr\hr}\tbr{a}ll	&	{\hr\hr}\tbr{a}r	&	{\hr\hr}\tbr{a}-ne	&	{\hr\hr}\tbr{e}my	&	{\hr\hr}\tbr{a}pn	&	{\hr\hr}\tbr{a}iy-ei\\
W. Marsela	&	{\hr\hr}\tbr{a}l	&	{\hr\hr}\tbr{a}r	&	{\hr\hr}\tbr{a}	&	{\hr\hr}\tbr{a}m	&	{\hr\hr}\tbr{a}pn	&	{\hr\hr}\tbr{a}y\\
E. Marsela	&	{\hr\hr}\tbr{a}l	&	{\hr\hr}\tbr{a}r	&	\hr\cg{(ulul)}	&	{\hr\hr}\tbr{a}m	&	{\hr\hr}\tbr{a}pen	&	{\hr\hr}\tbr{a}i\\
Serili	&	{\hr\hr}h\tbr{a}ll-e	&	{\hr\hr}h\tbr{a}r	&	{\hr\hr}\tbr{a}a-le	&	{\hr\hr}\tbr{a}m	&	{\hr\hr}\tbr{a}pel	&	{\hr\hr}\tbr{a}i-ɛ\\
SE. Babar	&	{\hr\hr}\tbr{a}l	&		&	{\hr\hr}\tbr{a}-u	&	{\hr\hr}\tbr{a}m	&	{\hr\hr}\tbr{a}pel	&	{\hr\hr}\tbr{a}i\\
Emplawas	&	{\hr\hr}\tbr{a}s-o	&	{\hr\hr}\tbr{a}r-e	&	{\hr\hr}\tbr{a}-lɛ	&	{\hr\hr}\tbr{a}m	&	{\hr\hr}\tbr{a}pl-i	&	{\hr\hr}\tbr{a}y-o\\
Imroing	&	{\hr\hr}\tbr{a}nn-e	&	{\hr\hr}\tbr{a}r-o	&	{\hr\hr}\tbr{a}i	&	\hr\cg{(iku)}	&	{\hr\hr}\tbr{a}pel-mu	&	{\hr\hr}\tbr{a}i-e\\
Tela	&	{\hr\hr}\tbr{a}nn-e	&	{\hr\hr}\tbr{a}r-e	&	{\hr\hr}\tbr{a}k-e	&	{\hr\hr}\tbr{a}mir	&	{\hr\hr}\tbr{a}pni	&	{\hr\hr}\tbr{a}i-ɛ\\
		\lspbottomrule
	\end{tabular}
		\begin{tablenotes}\tnsize
			\item[†]	{Cognates of **har ‘betel’ in \tsc{Peripheral Babar} languages
								are: Dawera \it{ar taw} (\it{taw} = ‘leaf’),
								Dai \it{har-on}, North Babar \it{ar kawe} (\it{kawe} = ‘leaf’),
								and Selaru \it{sar} ‘betel vine’.
								These words point to earlier **zar.}
			\item[‡]	{Cognates of **yai ‘boat’ include: Teun \it{yai},
								Nila (Wotay) \it{ai} [ʔai],
								 Nila (Kokroman) \it{yai},
								West Damer \it{osk-o}, Dawera \it{asy-ol}, Dai \it{ais-on}.
								Evidence for initial **y comes from retention of **a = \it{a},
								as well as the Teun and Nila cognates.
								For **y > {\0} compare *R > **y > {\0} in \trf{04-47}.}
		\end{tablenotes}
	\end{threeparttable}
%\end{adjustbox}
\end{table}

There are some morphemes which alternate between forms with /a/
and forms with a back rounded vowel.
Thus, for instance, \citet[98]{Steinhauer2009} gives the paradigm
of ‘die’ (from PMP *matay) in Southeast Babar as \tsc{1sg} \it{i-m\<y\>axy},
\tsc{2sg}/\tsc{2pl} \it{m-m\<y\>axy}, \tsc{3sg} \it{l-moxy},
\tsc{1pl.incl} \it{x-moxy}, \tsc{1pl.excl} \it{m-moxy}, \tsc{3pl} \it{t-moxy}.

Within \tsc{Nuclear South Babar},
all languages except West Marsela share **n > \it{l, n}.
Whether or not Tela belongs in this node is slightly ambiguous,
as it has undergone **l > \it{n} in nearly all cases.\footnote{
		There are also three sporadic instances of *l > \it{n} in \tsc{South
		Babar} languages other than Tela: *ma\tbr{l}ip > Serili \it{lye-mo\tbr{n}},
		Emplawas \it{lɛ-mu\tbr{n}}, Imroing \it{mo-mu\tbr{n}} ‘laugh’,
		*qa\tbr{l}ima > East Marsela \it{\tbr{n}im} ‘hand/arm’,
		and *da\tbr{l}əm > **ra\tbr{l}m > Serili \it{ram\tbr{n}-e}
		‘inside’.}
However, a handful of words have **n > \it{l}; *ə\tbr{n}əm > \it{wo\tbr{l}ɛm}
‘six’, *pa\tbr{n}iki > \it{\tbr{l}i-ɛ} ‘fruit bat’,
*mi\tbr{ñ}ak > \it{me\tbr{l}y-e}, and one word has *l = \it{l};
*\tbr{l}ima > \it{wu\tbr{l}im} ‘five’.
We propose that these words reflect the earlier change of **n/**l
> \it{l} and were not affected by later **l > \it{n}.\footnote{
		The latter
		change of Tela *l > \it{n} also affects loans,
		such as Malay \it{gula} > Tela \it{kun-e} ‘sugar’.
		This may indicate that it is a relatively recent change.}
These words thus provide evidence that Tela belongs in the node defined
by **n > **l, with subsequent Tela **l > \it{n} erasing evidence of this change
in most, but not all, words.

\citet{Zobel2018} shows that the change of **n > \it{l, n}
is conditioned.
**n is retained unchanged adjacent to PMP *a (which had shifted
to Proto-South Babar **ə in Zobel’s account) but shifts to \it{l} in other positions.
East Marsela and Serili have additional examples of **n = \it{n}
for which no consistent conditioning environment has been identified.
We cannot identify /u/ as the vowel adjacent to which **n = \it{n},
due to examples such as *bituqən > **kun > Emplawas \it{kul-o} ‘star’.
This means that *a > \it{u} must have occurred after **n > \it{l}.
Given that West Marsela has *a > \it{u} but does not usually
have **n > \it{l},\footnote{
		There are two examples of *n > \it{l}
		in available West Marsela data: *wa\tbr{n}an ‘right’ > West Marsela
		\it{wa\tbr{l}} (Central Marsela \it{wa\tbr{n}y},
		Serili \it{bya\tbr{l}el}) and *puqu\tbr{n} > West Marsela \it{u\tbr{l}}
		(Central Marsela \it{o\tbr{n}-e}, East Marsela \it{o\tbr{l}-e}, Serili \it{o\tbr{l}},
		Emplawas \it{o\tbr{s}y-o} (**olli), Imroing \it{o\tbr{l}y-e}).}
this implies that the change *a > \it{u} spread by diffusion
among \tsc{Nuclear South Babar} after **n > \it{l}.

Proto-South Babar **n is also usually retained unchanged before
/t/, and \citet[102]{Steinhauer2009} shows that for some morphemes
/l/ and /n/ synchronically alternate before /t/ in Southeast Babar.
Examples include \it{-tol} ‘ill’ \it{→ -ton{\tl}tol} ‘very ill’,
and \it{lewal} ‘language’ + \it{tool} ‘our’ → \it{lewan tool} ‘our language’.
Examples of 
**n > \it{l} are given in  \trf{04-58}
and
examples of
**n = \it{n} are given in \trf{04-60}.

\begin{table}\tcp{\tsc{Nuclear \tsc{South Babar}} **n > \it{l}}{04-58}
\begin{adjustbox}{width=\textwidth}
	\begin{threeparttable}
	{\st{2pt}\begin{tabular}{ll@{}l@{}llll@{}l}\lsptoprule
local gl.	&	{\hr}mouth	&	{\hr}snake	&	{\hr}tooth	&	{\hr}wound	&	{\hr}bat	&	{\hr}male	&	{\hr}full	\\
PMP	&	\hr*\tbb{ŋ}usuq	&	\hr*\tbb{n}ipay	&	\hr*\tbb{n}ipə\tbb{n}	&	\hr*(ma)-\tbb{n}uka	&	\hr*pa\tbb{n}iki	&	\hr*maRuqa\tbb{n}ay	&	\hr*pə\tbb{n}uq	\\
P.-Babar	&	**\tbb{n}uru\su{†}	&	**\tbb{n}ie	&	**\tbb{n}i\tbb{n}	&	**(m)\tbb{n}u	&	**\tbb{n}iki	&	**mayu\tbb{n}i	&	**pe\tbb{n}i	\\
P.-S. Babar	&	**\tbb{n}or	&	**\tbb{n}ie	&	**\tbb{n}i\tbb{n}	&	**(m)\tbb{n}o	&	**ni	&	**oi ma\tbb{n}i\su{‡}	&	**pe\tbb{n}i	\\ \midrule
C. Marsela	&	{\hr\hr}\tbb{n}or-ro	&	{\hr\hr}\tbb{n}i	&	{\hr\hr}\tbb{n}i\tbb{n}n	&	{\hr\hr}\tbb{n}o-ne	&	{\hr\hr}\tbb{n}i-ei	&	{\hr\hr}o myɛ\tbb{n}	&	{\hr\hr}\tbb{n}-pe\tbb{n}{\tl}pye\tbb{n}	\\
W. Marsela	&	{\hr\hr}\tbr{l}or	&	{\hr\hr}li\tbr{l}ikm	&	{\hr\hr}\tbb{n}i\tbb{n}	&	{\hr\hr}\tbb{n}o{\tl}no	&	{\hr\hr}\tbb{n}iy	&	{\hr\hr}oy mɛ\tbb{n}	&	{\hr\hr}\tbb{n}-pɛ\tbb{n}	\\
E. Marsela	&	{\hr\hr}\tbr{l}or	&	{\hr\hr}\tbr{l}e	&	{\hr\hr}\tbr{l}i\tbr{l}	&	{\hr\hr}\tbr{l}o	&	{\hr\hr}\tbr{l}i	&	{\hr\hr}oi mɛ\tbr{l}	&	{\hr\hr}pɛl{\tl}pɛ\tbr{l}	\\
Serili	&	{\hr\hr}\tbr{l}orh-e	&	{\hr\hr}\tbr{l}ɛ-ɛ	&	{\hr\hr}\tbr{l}i\tbr{l}	&	{\hr\hr}\tbr{l}o	&	{\hr\hr}\tbr{l}i-e	&	{\hr\hr}o myɛ\tbb{n}, omya\tbb{n}	&	{\hr\hr}pel{\tl}pye\tbr{l}	\\
SE. Babar	&	{\hr\hr}\tbr{l}or	&	{\hr\hr}\tbr{l}e	&	{\hr\hr}\tbr{l}i\tbr{l}	&	{\hr\hr}m\tbr{l}o	&	{\hr\hr}\tbr{l}i	&	{\hr\hr}oi mya\tbr{l}	&	{\hr\hr}l-pɛl{\tl}pi\tbr{l}	\\
Emplawas	&	{\hr\hr}\tbr{l}or-ɛ	&	{\hr\hr}li\tbr{l}ikm-o	&	{\hr\hr}\tbr{l}i\tbr{l}-wo	&	{\hr\hr}lɛ-\tbr{l}o	&	{\hr\hr}\tbr{l}i-o	&	{\hr\hr}pyɛ\tbr{l}y-o	&	{\hr\hr}pɛ{\tl}pyɛ\tbr{l}	\\
Imroing	&	{\hr\hr}\tbr{l}or-emu	&	{\hr\hr}i\tbr{l}ikəme	&	{\hr\hr}\tbr{l}i\tbr{l}-mu-te	&	{\hr\hr}lɛ-\tbr{l}o	&		&	{\hr\hr}o me\tbr{l}i	&	{\hr\hr}pɛl{\tl}pɛ\tbr{l}i-nili	\\
Tela	&	{\hr\hr}\tbb{n}ɔr-e	&	{\hr\hr}ne\tbb{n}ikmə	&	{\hr\hr}\tbb{n}i\tbb{n}-e	&	{\hr\hr}\tbb{n}ɛ-m\tbb{n}o	&	{\hr\hr}\tbr{l}i-ɛ	&	{\hr\hr}o me\tbb{n}-e	&	{\hr\hr}n-pen{\tl}piɛ\tbb{n}i	\\
\midrule\midrule
	\end{tabular}}
	{\st{2pt}\begin{tabular}{ll@{}ll@{}l@{}lll}
local gl.	&	{\hr}hear	&	{\hr}roast	&	{\hr}sand	&	{\hr}thousand	&	{\hr}kill	&	{\hr}fish	&	{\hr}star	\\
PMP	&	\hr*də\tbb{ŋ}əR	&	\hr*tu\tbb{n}u	&	\hr*qə\tbb{n}ay	&	\hr*Ribu	&	\hr*bu\tbb{n}uq	&	\hr*hika\tbb{n}	&	\hr*bituqə\tbb{n}	\\
P.-Babar	&		&	**tu\tbb{n}i	&	**e\tbb{n}i	&	**riwi\tbb{n}	&	**bu\tbb{n}i	&	**ika\tbb{n}	&	**tu\tbb{n}	\\
P-S. Babar	&	**re\tbb{n}	&	**ku\tbb{n}i	&	**i\tbb{n}i	&	**riwe\tbb{n}(i)	&	**bu\tbb{n}i	&	**ia\tbb{n}	&	**ku\tbb{n}	\\ \midrule
C. Marsela	&	{\hr\hr}rɛ\tbb{n}-ker	&	{\hr\hr}m-ku\tbb{n}-i	&	{\hr\hr}i\tbb{n}y	&	{\hr\hr}riw\tbb{n}i, riwe\tbb{n}	&	\hr\cg{(n-eni)}	&	{\hr\hr}i\tbb{n}-ei	&	{\hr\hr}ku\tbb{n}-ei	\\
W. Marsela	&	{\hr\hr}rɛ\tbb{n}	&	{\hr\hr}r-ku\tbb{n}-i	&	{\hr\hr}i\tbb{n}y	&	{\hr\hr}riw\tbb{n}	&	{\hr\hr}n-wu\tbb{n}	&	{\hr\hr}i\tbb{n}	&	{\hr\hr}ku\tbb{n}	\\
E. Marsela	&	{\hr\hr}m-re\tbr{l}	&	{\hr\hr}ku\tbb{n}-i	&	{\hr\hr}i\tbr{l}	&	{\hr\hr}rivi\tbr{l}i	&	{\hr\hr}m-pu\tbr{l}-i	&	{\hr\hr}ɛ\tbr{l}	&	{\hr\hr}ku\tbb{n}	\\
Serili	&	{\hr\hr}mə-ri\tbr{l}	&	{\hr\hr}tɛ-ku\tbr{l}	&	{\hr\hr}i\tbr{l}i-e	&	{\hr\hr}nibɛ\tbr{l}	&	{\hr\hr}lə-bu\tbr{l}	&	{\hr\hr}ɛ\tbr{l}	&	{\hr\hr}ku\tbr{l}, u\tbr{l}-e	\\
SE. Babar	&	{\hr\hr}ri\tbr{l}-ker	&	{\hr\hr}kuhu\tbr{l}	&	{\hr\hr}i\tbr{l}i	&	{\hr\hr}ribə\tbr{l}i	&	{\hr\hr}lɛ-bu\tbr{l}	&	{\hr\hr}ɛ\tbr{l}	&	{\hr\hr}ku\tbr{l}amol	\\
Emplawas	&	{\hr\hr}rɛ\tbr{l}(-ker)	&	{\hr\hr}ku\tbr{l}-i	&	{\hr\hr}i\tbr{l}y-o, ɛ\tbr{l}y-o	&	{\hr\hr}riw\tbr{l}	&	{\hr\hr}lə-wu\tbr{l}-i\su{Δ}	&	{\hr\hr}i\tbr{l}-o	&	{\hr\hr}ku\tbr{l}-o	\\
Imroing	&	\hr\cg{(nukainu)}	&	{\hr\hr}m-ku\tbr{l}-i	&	{\hr\hr}i\tbr{l}i-e	&	{\hr\hr}riuw\tbr{l}i	&	\hr\cg{(m-eli)}	&	{\hr\hr}i\tbr{l}-e	&	\hr\cg{(kowa{\tl}kowa)}	\\
Tela	&	\hr\cg{(nukain)}	&	{\hr\hr}me-kuk\tbb{n}-i	&	{\hr\hr}i\tbb{n}i-e	&	{\hr\hr}rive\tbb{n}i	&	\hr\cg{(n-enɛ)}	&	{\hr\hr}ii\tbb{n}-e	&	{\hr\hr}ku\tbb{n}-e	\\
		\lspbottomrule
	\end{tabular}}
		\begin{tablenotes}\tnsize
			\item[†] {PMP *ŋusuq has lexically specific irregular *s (> **d) > **r
								in all known \tsc{Southwest Maluku} languages.}
			\item[‡]	{Reflexes of *maRuqanay `man, male' are compounded
									with of **oi `person'. The final vowel of this stem
									often undergoes metathesis with the following morpheme.}
			\item[Δ] {Emplawas \it{lə-wul-i} ‘kill’ is from \citet[433]{Taber1993}.
								\citet{Edwards2023} has \it{l-wal} ‘hit, kill’.}
		\end{tablenotes}
	\end{threeparttable}
\end{adjustbox}
\end{table}

\begin{table}\tcp{\tsc{South Babar} **n = \it{n}}{04-60}
%\begin{adjustbox}{width=\textwidth}
	\st{3pt}\begin{tabular}{ll@{}l@{}l@{}lll}\lsptoprule
local gl.	&	{\hr}child	&	{\hr}bird	&	{\hr}name	&	{\hr}wind	&	{\hr}swim	&	{\hr}tongue\\
PMP	&	\hr*a\tbb{n}ak	&	\hr*ma\tbb{n}uk	&	\hr*ŋaja\tbb{n}	&	\hr*haŋi\tbb{n}	&	\hr*\tbb{n}aŋuy	&	\\
P.-Babar	&	**a\tbb{n}	&	**ma\tbb{n}u	&	**\tbb{n}a\tbb{n}	&	**aŋi\tbb{n}	&	** we\tbb{n}a\tbb{ŋ}ui	&	**\tbb{n}ama\\
P.-S. Babar	&	**a\tbb{n}	&	**ma\tbb{n}i(an)	&	**\tbb{n}a\tbb{n}	&	**a\tbb{nn}i	&	**we\tbb{n}a\tbb{n}i	&	**\tbb{n}am\\ \midrule
C. Marsela	&	{\hr\hr}ɛ\tbb{n}-e	&	{\hr\hr}mɛ\tbb{n}yɛn	&	{\hr\hr}\tbb{n}ɛ\tbb{n}	&	{\hr\hr}ɛ\tbb{nn}on	&	{\hr\hr}ne-wyɛ\tbb{n}	&	{\hr\hr}\tbb{n}ɛm\\
W. Marsela	&	{\hr\hr}u\tbb{n}	&	{\hr\hr}mu\tbb{n}y	&	{\hr\hr}\tbb{n}u\tbb{n}	&	{\hr\hr}u\tbb{nn}	&	{\hr\hr}mo-wnyɛ\tbb{n}	&	{\hr\hr}\tbb{n}um\\
E. Marsela	&	{\hr\hr}u\tbb{n}	&	{\hr\hr}mo\tbb{n}yen	&	{\hr\hr}\tbb{n}o\tbb{n}	&	{\hr\hr}u\tbb{n}	&	{\hr\hr}mo-w\tbb{n}o\tbb{n}	&	{\hr\hr}\tbb{n}om\\
Serili	&	{\hr\hr}u\tbb{n}-e	&	{\hr\hr}mo\tbb{n}ɛ	&	{\hr\hr}\tbb{n}o\tbb{n}	&	{\hr\hr}u\tbb{n}e\tbb{n}	&	{\hr\hr}lyɛ-mo\tbb{n}o\tbb{n}	&	{\hr\hr}\tbb{n}om\\
SE. Babar	&	{\hr\hr}u\tbb{n}	&	{\hr\hr}mo\tbb{n}yan	&	{\hr\hr}\tbb{n}o\tbb{n}	&	{\hr\hr}u\tbb{n}	&	{\hr\hr}ke-be\tbb{n}o\tbb{n}	&	{\hr\hr}\tbb{n}om\\
Emplawas	&	{\hr\hr}u\tbb{n}-ɛ	&	{\hr\hr}mu\tbb{n}yɛl-o	&	{\hr\hr}\tbb{n}u\tbb{n}-o	&	{\hr\hr}usy-o	&	{\hr\hr}lɛ-w\tbb{n}u\tbb{n}	&	{\hr\hr}\tbb{n}up-ɛ\\
Imroing	&	{\hr\hr}u\tbb{n}-ɛ-mu	&	{\hr\hr}mu\tbb{n}e	&	{\hr\hr}\tbb{n}u\tbb{n}-mu	&	{\hr\hr}u\tbb{nn}-ɛ	&	{\hr\hr}mo-u\tbr{l}u\tbr{l}i	&	{\hr\hr}\tbb{n}om-ne\\
Tela	&	{\hr\hr}u\tbb{n}-e	&	{\hr\hr}mu\tbb{n}e	&	{\hr\hr}\tbb{n}u\tbb{n}-mu	&	{\hr\hr}u\tbb{nn}-e	&	{\hr\hr}ne-w\tbb{n}u\tbb{n}i	&	{\hr\hr}\tbb{n}um-mu\\
		\lspbottomrule
	\end{tabular}
%\end{adjustbox}
\end{table}

Emplawas has three unusual sound changes involving **l: **ll > \it{s},
**lm/ml > \it{p}, and **rl > \it{tt}.
Examples of **ll > \it{s} include: *bulan > **wu\tbr{ll} > \it{wo\tbr{s}-o}
‘moon, month’, *kawanan > **wa\tbr{ll} > \it{wa\tbr{s}} ‘right (side)’,
*tələn > **ke\tbr{ll} > \it{l-ke\tbr{s}} ‘s/he swallows’,
and *zalan > **a\tbr{ll} > \it{a\tbr{s}-o} ‘path, road’.
An example of *ml > \it{p} is *ma-nuka > **\tbr{ml}o > \it{\tbr{p}o} ‘wound’.
Two examples of *rl > \it{t(t)} are: *ijuŋ > **(h)i\tbr{rl} > \it{(h)i\tbr{tt}-o}
‘my nose’ and *kudən > **o\tbr{rl} > \it{o\tbr{tt}-o} ‘pot’.
These changes result in morphophonemic changes in Emplawas,
shown by the inflected forms in \trf{EmpPosPar}.

\begin{table}\caption[Emplawas possessive paradigms]{Emplawas possessive paradigms\su{†}}\label{tab:EmpPosPar}
\begin{adjustbox}{width=\textwidth}
	\begin{threeparttable}\st{2pt}
	\begin{tabular}{ll|l@{}l@{}l@{}l@{ }l}\lsptoprule
local gl.	&	\tsc{gen}	&	{\hr}eye	&	{\hr}tooth	&	{\hr}ear		&	{\hr}arm/hand	&	{\hr}house	\\	
PMP	&		&	\hr*mata	&	\hr*nipən	&	\hr*taliŋa			&	\hr*qalima	&	\hr*Rumaq	\\	
P.-S. Babar&		&	**mak	&	**lil	&	**klil			&	**lim	&	***yum	\\	\midrule
\tsc{1sg}	&	{-(w)o}	&	\hr\hr{muk(w)o}	&	\hr\hr{lilwo}	&	\hr\hr{klilwo}		&	\hr\hr{limo}	&	\hr\hr{imwo}	\\	
\tsc{2sg}	&	{-mo}	&	\hr\hr{mukmo}	&	\hr\hr{li\tbr{pp}o}	&	\hr\hr{kli\tbr{pp}o}		&	\hr\hr{li\tbr{pp}o}	&	\hr\hr{ippo}	\\	
\tsc{3sg}	&	{-lɛ/-lyo}	&	\hr\hr{muklɛ/muklyo}	&	\hr\hr{li\tbr{s}ɛ/li\tbr{s}yo}	&	\hr\hr{kli\tbr{s}ɛ/kli\tbr{s}yo}		&	\hr\hr{li\tbr{pp}ɛ}	&	\hr\hr{i\tbr{p}yo/i\tbr{pp}ɛ}	\\	
\tsc{1pl.incl}\su{‡}	&	{-tɛ}	&	\hr\hr{muktɛtɛr}	&	\hr\hr{litɛtɛr}	&	\hr\hr{klitɛtɛr}		&	\hr\hr{limtɛtɛr}	&	\hr\hr{imtɛtɛr}	\\	
%\tsc{1pl.excl}	&	{-mumyo}	&		&		&			&		&	\hr\hr{im-mumyo}	\\	
\tsc{2pl}	&	{-mi/-m(y)o}	&	\hr\hr{mukmitɛr}	&	\hr\hr{li\tbr{p}otɛr}	&	\hr\hr{kli\tbr{p}(y)otɛr}		&	\hr\hr{li\tbr{pp}itɛr}	&	\hr\hr{}	\\	
\tsc{3pl}	&	{-lyo}	&	\hr\hr{muklyotɛr}	&	\hr\hr{li\tbr{s}yotɛr}	&	\hr\hr{kli\tbr{s}yotɛr}		&	\hr\hr{li\tbr{p}yotɛr}	&	\hr\hr{i\tbr{p}yo}	\\	
		\lspbottomrule
	\end{tabular}
		\begin{tablenotes}\tnsize
			\item[†] {\it{-tɛr}, which occurs on several forms, is a plural suffix.}
			\item[‡] {\tsc{1pl.excl} forms are poorly attested.
								For these nouns, only \it{im-umyo/im-imyo} `our house'
								has been encountered.
								Other nouns have \tsc{1pl.excl} \it{-mumyo/-mumyɛ}.}
		\end{tablenotes}
	\end{threeparttable}
\end{adjustbox}
\end{table}
